% Options for packages loaded elsewhere
\PassOptionsToPackage{unicode}{hyperref}
\PassOptionsToPackage{hyphens}{url}
\documentclass[
  russian,
]{report}
\usepackage{xcolor}
\usepackage{amsmath,amssymb}
\setcounter{secnumdepth}{-\maxdimen} % remove section numbering
\usepackage{iftex}
\ifPDFTeX
  \usepackage[T1]{fontenc}
  \usepackage[utf8]{inputenc}
  \usepackage{textcomp} % provide euro and other symbols
\else % if luatex or xetex
  \usepackage{unicode-math} % this also loads fontspec
  \defaultfontfeatures{Scale=MatchLowercase}
  \defaultfontfeatures[\rmfamily]{Ligatures=TeX,Scale=1}
\fi
\usepackage{lmodern}
\ifPDFTeX\else
  % xetex/luatex font selection
\fi
% Use upquote if available, for straight quotes in verbatim environments
\IfFileExists{upquote.sty}{\usepackage{upquote}}{}
\IfFileExists{microtype.sty}{% use microtype if available
  \usepackage[]{microtype}
  \UseMicrotypeSet[protrusion]{basicmath} % disable protrusion for tt fonts
}{}
\makeatletter
\@ifundefined{KOMAClassName}{% if non-KOMA class
  \IfFileExists{parskip.sty}{%
    \usepackage{parskip}
  }{% else
    \setlength{\parindent}{0pt}
    \setlength{\parskip}{6pt plus 2pt minus 1pt}}
}{% if KOMA class
  \KOMAoptions{parskip=half}}
\makeatother
\ifLuaTeX
\usepackage[bidi=basic,shorthands=off,provide=*]{babel}
\else
\usepackage[bidi=default,shorthands=off,provide=*]{babel}
\fi
\ifLuaTeX
  \usepackage{selnolig} % disable illegal ligatures
\fi
\setlength{\emergencystretch}{3em} % prevent overfull lines
\providecommand{\tightlist}{%
  \setlength{\itemsep}{0pt}\setlength{\parskip}{0pt}}
\usepackage{bookmark}
\IfFileExists{xurl.sty}{\usepackage{xurl}}{} % add URL line breaks if available
\urlstyle{same}
\hypersetup{
  pdflang={ru},
  hidelinks,
  pdfcreator={LaTeX via pandoc}}

\title{Skill Mismatch и Принцип Питера\\
Социологический анализ проблемы несоответствия компетенций в найме и
интеллектуальная система её решения}
\author{Московский авиационный институт\\
(национальный исследовательский университет)\\
Институт №8 «Фундаментальная информатика и информационные технологии»\\
\strut \\
Авторы: студенты кафедры 316}
\date{Москва --- 2025}

\begin{document}
\maketitle

{
\setcounter{tocdepth}{2}
\tableofcontents
}
\chapter{1. Введение}\label{ux432ux432ux435ux434ux435ux43dux438ux435}

\section{1.1. Актуальность
исследования}\label{ux430ux43aux442ux443ux430ux43bux44cux43dux43eux441ux442ux44c-ux438ux441ux441ux43bux435ux434ux43eux432ux430ux43dux438ux44f}

В условиях стремительной цифровой трансформации мировой экономики
проблема \textbf{несоответствия компетенций} (Skill Mismatch) становится
одним из ключевых барьеров для устойчивого экономического роста. По
данным Международной организации труда (ILO), от 25 до 45\% работников в
развитых странах занимают должности, не соответствующие их реальным
навыкам --- они либо переквалифицированы, либо недоквалифицированы
относительно требований рабочего места.

Параллельно с этим действует \textbf{Принцип Питера} --- организационный
феномен, при котором сотрудники продвигаются по карьерной лестнице до
тех пор, пока не достигнут уровня собственной некомпетентности. Этот
эффект системно усиливает skill mismatch внутри организаций: вчерашний
эффективный программист становится неэффективным менеджером, а его
позицию занимает сотрудник, ещё не готовый к задачам такого уровня.

Совокупность этих явлений порождает \textbf{парадокс современного
найма}: высокооплачиваемые специалисты вносят непропорционально малый
вклад в конечный продукт, тогда как низкооплачиваемые работники
критически важны для операционной деятельности. Традиционные системы
оценки персонала, основанные на субъективных интервью, самоотчётах и
формальных аттестациях, не способны объективно выявить и исправить эти
диспропорции.

\section{1.2. Цель и
задачи}\label{ux446ux435ux43bux44c-ux438-ux437ux430ux434ux430ux447ux438}

\textbf{Цель работы} --- исследовать социологические и экономические
аспекты проблемы Skill Mismatch и Принципа Питера, а также предложить
интеллектуальную систему объективной оценки компетенций на основе
анализа реальных рабочих артефактов.

\textbf{Задачи:}

\begin{enumerate}
\def\labelenumi{\arabic{enumi}.}
\tightlist
\item
  Определить сущность и масштаб проблемы Skill Mismatch на основе данных
  ILO, OECD и академических исследований
\item
  Проанализировать механизм действия Принципа Питера и его влияние на
  организационную эффективность
\item
  Оценить экономические потери от неоптимального распределения
  человеческого капитала
\item
  Предложить математическую модель объективной оценки компетенций
  специалистов
\item
  Описать платформу DevMetrics как практическую реализацию предложенного
  подхода
\end{enumerate}

\section{1.3. Объект и предмет
исследования}\label{ux43eux431ux44aux435ux43aux442-ux438-ux43fux440ux435ux434ux43cux435ux442-ux438ux441ux441ux43bux435ux434ux43eux432ux430ux43dux438ux44f}

\begin{itemize}
\tightlist
\item
  \textbf{Объект} --- процессы найма, оценки и продвижения специалистов
  в ИТ-индустрии
\item
  \textbf{Предмет} --- закономерности формирования skill mismatch и
  методы его выявления на основе автоматизированного анализа рабочих
  артефактов
\end{itemize}

\section{1.4.
Методология}\label{ux43cux435ux442ux43eux434ux43eux43bux43eux433ux438ux44f}

Исследование основывается на:

\begin{itemize}
\tightlist
\item
  \textbf{Статистическом анализе} данных ILO и OECD по 37 странам
\item
  \textbf{Социологическом анализе} явления Принципа Питера
\item
  \textbf{Математическом моделировании} системы грейдирования с функцией
  затухания Грина
\item
  \textbf{Программной реализации} платформы DevMetrics с применением ИИ
  (LLM) для анализа git-коммитов и метрик кода
\end{itemize}

\newpage

\end{document}
